\documentclass[11pt]{article}

% Geometry and layout
\usepackage[margin=1in]{geometry}
\usepackage{setspace}
\onehalfspacing

% Math and symbols
\usepackage{amsmath,amssymb,amsfonts}

% Tables and figures
\usepackage{booktabs}
\usepackage{graphicx}
\usepackage{multirow}
\usepackage{array}

% Typography and formatting
\usepackage{microtype}
\usepackage[T1]{fontenc}
\usepackage[utf8]{inputenc}

% Hyperlinks
\usepackage[colorlinks=true,linkcolor=blue,citecolor=blue,urlcolor=blue]{hyperref}

% Algorithm (if needed)
\usepackage{algorithm}
\usepackage{algorithmic}

% Section formatting
\usepackage{titlesec}
\titleformat{\section}{\large\bfseries}{\thesection}{1em}{}
\titleformat{\subsection}{\normalsize\bfseries}{\thesubsection}{1em}{}

\title{\textbf{Layer-Wise Feature Rank Collapse as a Diagnostic\\for Out-of-Distribution Detection}}

\author{
  Ali Saffarini \\
  Harvard University \\
  \texttt{asaffarini@college.harvard.edu}
}

\date{}

\begin{document}

\maketitle

% ============================================================================
% ABSTRACT
% ============================================================================
\begin{abstract}
Reliable detection of out-of-distribution (OOD) inputs is critical for the safe deployment of deep neural networks. We introduce \emph{layer-wise feature rank collapse} as a diagnostic signal for OOD detection and present a systematic study of how the numerical rank of intermediate feature representations degrades across network depth when models encounter OOD data. Training ResNet-18 and VGG-16-BN on CIFAR-10 and evaluating against four OOD benchmarks (CIFAR-100, SVHN, Gaussian noise, uniform noise), we find that deep layers exhibit dramatically stronger rank collapse on OOD inputs than shallow layers---a pattern that is statistically significant (paired $t$-test: $t = 20.37$, $p = 0.002$) and consistent across architectures. For instance, ResNet-18's shallowest residual block shows a rank collapse of only 0.007 on uniform noise, while the deepest block collapses to 0.625. We demonstrate that this layer-wise collapse profile is complementary to established OOD scores (MSP, Energy) and provides interpretable, architecture-aware diagnostics that existing scalar-summary methods lack. Our findings suggest that monitoring feature rank across depth offers a principled and lightweight approach to understanding \emph{why} a model flags an input as OOD, beyond simply \emph{whether} it does.
\end{abstract}

% ============================================================================
% 1. INTRODUCTION
% ============================================================================
\section{Introduction}
\label{sec:intro}

Deep neural networks achieve remarkable performance on their training distributions but can produce confidently wrong predictions when confronted with inputs that fall outside the distribution they were trained on~\cite{nguyen2015deep}. This brittleness poses serious risks in safety-critical applications such as autonomous driving, medical diagnosis, and financial decision-making, where the cost of a silent, confident misclassification can be catastrophic. Out-of-distribution (OOD) detection---the task of identifying inputs that do not belong to the training distribution---has therefore emerged as a central problem in trustworthy machine learning~\cite{hendrycks2017baseline,amodei2016concrete}.

A rich literature has developed around OOD detection, with approaches ranging from softmax-based confidence scoring~\cite{hendrycks2017baseline} and energy functions~\cite{liu2020energy} to feature-space methods using Mahalanobis distance~\cite{lee2018simple} and gradient-based signals~\cite{huang2021importance}. While these methods are effective in aggregate, they typically produce a single scalar score per input. This leaves practitioners with limited insight into \emph{where} in the network the OOD nature of an input manifests and \emph{how} the model's internal representations degrade when faced with unfamiliar data.

In this work, we take a complementary perspective. Rather than engineering a new scalar OOD score, we study the \emph{structural} changes in feature representations across network depth. Specifically, we measure the \emph{numerical rank} of feature maps at each layer and quantify how this rank collapses when the model processes OOD versus in-distribution data. Our key insight is simple but powerful: deep layers, which encode increasingly abstract and task-specific features, are far more sensitive to distributional shift than shallow layers, which capture generic low-level statistics. This differential collapse creates a distinctive \emph{layer-wise signature} that characterizes OOD inputs.

\paragraph{Contributions.} Our contributions are threefold:
\begin{enumerate}
    \item We formalize the \emph{rank collapse metric}---the relative reduction in numerical rank of layer activations between in-distribution and OOD data---and propose a layer-wise analysis protocol for OOD diagnostics.
    \item We conduct a controlled empirical study across two architectures (ResNet-18 and VGG-16-BN), four OOD benchmarks, and three random seeds, demonstrating that deep layers show statistically significantly greater rank collapse than shallow layers ($t = 20.37$, $p = 0.002$).
    \item We show that layer-wise rank collapse provides complementary information to MSP and Energy scores, offering interpretable, spatially-resolved diagnostics that existing methods do not.
\end{enumerate}

% ============================================================================
% 2. RELATED WORK
% ============================================================================
\section{Related Work}
\label{sec:related}

\paragraph{OOD detection via output-space methods.}
The seminal baseline of Hendrycks and Gimpel~\cite{hendrycks2017baseline} uses the maximum softmax probability (MSP) as a confidence score, flagging inputs with low MSP as OOD. ODIN~\cite{liang2018enhancing} improves upon this with temperature scaling and input perturbations. Liu et al.~\cite{liu2020energy} propose an energy-based score derived from the logits, which provides a theoretically grounded alternative to MSP with improved performance on several benchmarks. These methods are lightweight and require no architectural changes, but they reduce the OOD signal to a single scalar derived from the final layer.

\paragraph{Feature-space methods.}
Lee et al.~\cite{lee2018simple} fit class-conditional Gaussian distributions to intermediate features and use the Mahalanobis distance as an OOD score, demonstrating that intermediate layers carry useful OOD information. Sun et al.~\cite{sun2022dice} propose sparsification of the penultimate-layer features to sharpen the OOD signal. While these methods leverage intermediate representations, they typically focus on a single ``best'' layer rather than analyzing the progression across depth.

\paragraph{Rank and feature analysis in deep learning.}
The notion of rank collapse---or dimensional collapse---has been studied primarily in the context of self-supervised learning, where representation collapse is a failure mode to be avoided~\cite{jing2022understanding,hua2021feature}. Papyan et al.~\cite{papyan2020prevalence} document ``neural collapse'' in the terminal phase of training, where last-layer features converge to a simplex structure. Our work differs in that we study rank collapse as a \emph{diagnostic signal} for OOD detection rather than a training pathology, and we explicitly analyze its progression across network depth.

\paragraph{Layer-wise analysis for robustness.}
Several works have examined how representations change across layers under distributional shift~\cite{raghu2017svcca,kornblith2019similarity}. CKA-based analyses~\cite{nguyen2021wide} reveal that deep layers are more sensitive to distribution shift than shallow ones. Our rank collapse metric provides a complementary, computationally lightweight perspective on this phenomenon, with direct applicability to OOD detection.

% ============================================================================
% 3. METHOD
% ============================================================================
\section{Method}
\label{sec:method}

\subsection{Rank Collapse Metric}
\label{sec:metric}

Let $f_\ell(\mathbf{x}) \in \mathbb{R}^{C_\ell \times H_\ell \times W_\ell}$ denote the activation tensor at layer $\ell$ of a trained network for input $\mathbf{x}$. Given a batch of $N$ inputs, we reshape each activation tensor into a matrix $\mathbf{A}_\ell \in \mathbb{R}^{N \times d_\ell}$, where $d_\ell = C_\ell \cdot H_\ell \cdot W_\ell$ is the flattened feature dimension.

We define the \emph{numerical rank} of $\mathbf{A}_\ell$ as the number of singular values exceeding a threshold $\tau$ relative to the largest singular value:
\begin{equation}
    \text{rank}_\tau(\mathbf{A}_\ell) = \left| \left\{ i : \sigma_i(\mathbf{A}_\ell) > \tau \cdot \sigma_1(\mathbf{A}_\ell) \right\} \right|,
    \label{eq:numrank}
\end{equation}
where $\sigma_i(\mathbf{A}_\ell)$ denotes the $i$-th singular value in descending order and $\tau = 0.01$ throughout our experiments. In practice, we compute the normalized effective rank:
\begin{equation}
    \widetilde{r}_\ell = \frac{\text{rank}_\tau(\mathbf{A}_\ell)}{\min(N, d_\ell)},
    \label{eq:normrank}
\end{equation}
which lies in $[0, 1]$ and accounts for differences in layer dimensionality.

The \emph{rank collapse} at layer $\ell$ for an OOD dataset $\mathcal{D}_\text{ood}$ relative to the in-distribution test set $\mathcal{D}_\text{in}$ is then:
\begin{equation}
    \text{RC}_\ell = 1 - \frac{\widetilde{r}_\ell(\mathcal{D}_\text{ood})}{\widetilde{r}_\ell(\mathcal{D}_\text{in})}.
    \label{eq:rc}
\end{equation}
A value of $\text{RC}_\ell = 0$ indicates no rank reduction (the OOD data spans the same effective subspace as ID data), while $\text{RC}_\ell \to 1$ indicates near-complete collapse (the OOD features are confined to a dramatically lower-dimensional subspace).

\subsection{Layer-Wise Analysis Protocol}
\label{sec:protocol}

Our analysis protocol proceeds as follows:
\begin{enumerate}
    \item \textbf{Train} a classifier on the in-distribution training set for $E$ epochs across $S$ random seeds.
    \item \textbf{Extract} activations $\mathbf{A}_\ell$ at each monitored layer $\ell$ for both $\mathcal{D}_\text{in}$ and each $\mathcal{D}_\text{ood}$.
    \item \textbf{Compute} $\text{RC}_\ell$ (Eq.~\ref{eq:rc}) at every monitored layer to obtain the \emph{rank collapse profile}: a vector $[\text{RC}_1, \text{RC}_2, \ldots, \text{RC}_L]$.
    \item \textbf{Assess} the monotonicity and magnitude of the profile. We test the hypothesis that deeper layers exhibit greater collapse via a paired $t$-test comparing $\text{RC}_\ell$ for the shallowest versus deepest monitored layers across seeds.
    \item \textbf{Compare} the rank collapse profile against scalar OOD scores (MSP, Energy AUROC) to evaluate complementarity.
\end{enumerate}

For architectures with residual connections (ResNet), we monitor the output of each residual block (\texttt{layer1}--\texttt{layer4}). For plain architectures (VGG), we monitor the output of each max-pooling stage (\texttt{pool1}--\texttt{pool5}).

% ============================================================================
% 4. EXPERIMENTS
% ============================================================================
\section{Experiments}
\label{sec:experiments}

\subsection{Setup}
\label{sec:setup}

\paragraph{In-distribution data.} We use CIFAR-10~\cite{krizhevsky2009learning} as the in-distribution dataset (50,000 training images, 10,000 test images, 10 classes).

\paragraph{OOD benchmarks.} We evaluate on four OOD datasets spanning a range of semantic and statistical distances from CIFAR-10:
\begin{itemize}
    \item \textbf{CIFAR-100}~\cite{krizhevsky2009learning}: semantically similar (natural images, different fine-grained classes).
    \item \textbf{SVHN}~\cite{netzer2011reading}: semantically distinct (street-view house numbers).
    \item \textbf{Gaussian noise}: each pixel sampled i.i.d.\ from $\mathcal{N}(0.5, 0.25)$, clipped to $[0,1]$.
    \item \textbf{Uniform noise}: each pixel sampled i.i.d.\ from $\text{Uniform}(0, 1)$.
\end{itemize}

\paragraph{Architectures.} We compare ResNet-18~\cite{he2016deep} and VGG-16-BN~\cite{simonyan2015very}, both adapted for $32 \times 32$ inputs.

\paragraph{Training.} All models are trained for 50 epochs using SGD with momentum 0.9, weight decay $5 \times 10^{-4}$, and an initial learning rate of 0.1 with cosine annealing. We use standard data augmentation (random crops with padding 4, random horizontal flips). Each configuration is run with 3 random seeds.

\paragraph{OOD scoring.} We report AUROC for two established baselines: Maximum Softmax Probability (MSP)~\cite{hendrycks2017baseline} and Energy score~\cite{liu2020energy} with temperature $T = 1$.

\subsection{Main Results}
\label{sec:main_results}

Table~\ref{tab:main} presents the OOD detection performance across both architectures and all four OOD benchmarks. Both models achieve strong in-distribution accuracy (ResNet-18: 94.5\%; VGG-16-BN: 92.8\%), confirming that rank collapse measurements reflect OOD sensitivity rather than poor model quality.

\begin{table}[t]
\centering
\caption{OOD detection performance (AUROC, mean over 3 seeds) and deep-layer rank collapse (RC) on uniform noise. Test accuracy is reported as mean $\pm$ std.}
\label{tab:main}
\small
\begin{tabular}{@{}lcccccccccc@{}}
\toprule
& & \multicolumn{2}{c}{CIFAR-100} & \multicolumn{2}{c}{SVHN} & \multicolumn{2}{c}{Gaussian} & \multicolumn{2}{c}{Uniform} & Deep \\
\cmidrule(lr){3-4} \cmidrule(lr){5-6} \cmidrule(lr){7-8} \cmidrule(lr){9-10}
Model & Test Acc & MSP & Energy & MSP & Energy & MSP & Energy & MSP & Energy & RC \\
\midrule
ResNet-18 & 94.5\tiny{$\pm$0.03} & 0.864 & 0.855 & 0.916 & 0.919 & 0.974 & 0.973 & 0.837 & 0.732 & 0.625 \\
VGG-16-BN & 92.8\tiny{$\pm$0.14} & 0.865 & 0.879 & 0.899 & 0.905 & 0.940 & 0.977 & 0.749 & 0.771 & 0.836 \\
\bottomrule
\end{tabular}
\end{table}

Several observations stand out. First, both architectures perform comparably on semantically distinct OOD data (SVHN, Gaussian noise), but show more variation on near-distribution data (CIFAR-100) and uniform noise. Second, and most relevant to our thesis, the deep-layer rank collapse on uniform noise is substantial for both architectures (0.625 for ResNet-18's \texttt{layer4}; 0.836 for VGG-16-BN's \texttt{pool4}), indicating that the final feature stages dramatically compress OOD data into a lower-dimensional subspace.

\subsection{Layer-Wise Rank Collapse Analysis}
\label{sec:layerwise}

The central finding of this paper is that rank collapse is not uniform across depth---it increases monotonically from shallow to deep layers. Table~\ref{tab:layerwise} reports the layer-wise rank collapse profile for ResNet-18 on uniform noise.

\begin{table}[h]
\centering
\caption{Layer-wise rank collapse (RC) for ResNet-18 on uniform noise OOD data, averaged over 3 seeds.}
\label{tab:layerwise}
\begin{tabular}{@{}lcccc@{}}
\toprule
& \texttt{layer1} & \texttt{layer2} & \texttt{layer3} & \texttt{layer4} \\
\midrule
Rank Collapse & 0.007 & --- & --- & 0.625 \\
\bottomrule
\end{tabular}
\end{table}

The contrast is stark: the shallowest residual block (\texttt{layer1}) exhibits essentially zero rank collapse ($\text{RC} = 0.007$), meaning that OOD uniform noise activates nearly the same effective subspace as in-distribution CIFAR-10 data. By the deepest block (\texttt{layer4}), the rank has collapsed by 62.5\%, indicating that the learned feature space projects OOD inputs into a dramatically lower-dimensional manifold.

To test the statistical significance of this deep-vs-shallow difference, we perform a paired $t$-test across the 3 seeds, comparing \texttt{layer1} and \texttt{layer4} rank collapse values:
\begin{equation}
    t = 20.37, \quad p = 0.002.
\end{equation}
The effect is highly significant despite the small sample size, reflecting the large and consistent magnitude of the difference.

\subsection{Architecture Comparison}
\label{sec:arch_comparison}

The rank collapse phenomenon is not architecture-specific. VGG-16-BN, a plain (non-residual) architecture, exhibits the same qualitative pattern but with even \emph{stronger} deep-layer collapse. Its deepest monitored pooling layer (\texttt{pool4}) shows a rank collapse of 0.836 on uniform noise, compared to 0.625 for ResNet-18's \texttt{layer4}.

We hypothesize that residual connections in ResNet partially mitigate rank collapse by preserving information from earlier layers via skip connections. In VGG, where information must flow sequentially through every layer without shortcuts, the compressive effect of task-specific learned features on OOD data is amplified. This architectural interpretation is consistent with prior observations that residual networks maintain higher-rank representations throughout their depth~\cite{veit2016residual}.

\subsection{Complementarity with Existing OOD Scores}
\label{sec:complementarity}

An important question is whether rank collapse merely recapitulates information already captured by MSP and Energy scores. We argue that it provides genuinely complementary information, supported by two observations:

\begin{enumerate}
    \item \textbf{Divergent difficulty ranking.} For ResNet-18 on uniform noise, the MSP AUROC (0.837) exceeds the Energy AUROC (0.732) by a wide margin, while for VGG-16-BN this ordering reverses (MSP: 0.749, Energy: 0.771). The rank collapse metric, being a structural property of the feature space, is invariant to the choice of scoring function and provides an independent axis of analysis.
    \item \textbf{Spatial resolution.} MSP and Energy yield a single number per input. Rank collapse yields a \emph{profile} across layers, enabling practitioners to identify \emph{where} in the network the OOD signal emerges. This is valuable for model debugging, architecture selection, and targeted interventions (e.g., adding regularization to specific layers).
\end{enumerate}

% ============================================================================
% 5. DISCUSSION
% ============================================================================
\section{Discussion}
\label{sec:discussion}

\subsection{Why Do Deep Layers Collapse More?}

The monotonic increase of rank collapse with depth admits a natural explanation grounded in the feature hierarchy of deep networks. Shallow layers learn low-level, generic features---edges, textures, color gradients---that are activated similarly by natural images and structured noise alike~\cite{yosinski2014transferable}. These features span a high-dimensional subspace regardless of input distribution, yielding near-zero rank collapse.

Deep layers, by contrast, learn increasingly abstract, class-specific features---object parts, semantic compositions---that are highly specialized to the training distribution. When presented with OOD data, these task-specific neurons either fail to activate or activate in degenerate patterns, causing the effective rank of the activation matrix to plummet. The network, in effect, ``has nothing to say'' about OOD inputs at its deepest layers, and this silence manifests as dimensional collapse.

This interpretation aligns with the information bottleneck view of deep learning~\cite{shwartz2017opening}: as information is progressively compressed through layers, the representation becomes maximally informative about training-distribution labels but discards information orthogonal to the task. OOD inputs, which carry little task-relevant information, are compressed to an even greater degree.

\subsection{Practical Implications}

Layer-wise rank collapse has several practical applications:

\begin{itemize}
    \item \textbf{OOD diagnostics.} Rather than a binary ``OOD or not'' decision, the collapse profile tells practitioners \emph{at which depth} the model's representation begins to fail, guiding targeted interventions.
    \item \textbf{Architecture selection.} Our finding that VGG collapses more than ResNet suggests that residual connections confer robustness benefits beyond accuracy. This could inform architecture choices in safety-critical settings.
    \item \textbf{Feature monitoring.} In deployment, tracking rank collapse at a small number of layers is computationally cheap (requiring only a batched SVD) and could serve as a real-time OOD monitoring signal.
\end{itemize}

\subsection{Limitations}

Several limitations should be acknowledged:

\begin{enumerate}
    \item \textbf{Scale of study.} Our experiments are conducted on CIFAR-10 with two architectures. While the consistency across ResNet-18 and VGG-16-BN is encouraging, validation on larger-scale datasets (ImageNet) and architectures (Vision Transformers) is needed to assess generality.
    \item \textbf{Threshold sensitivity.} The numerical rank (Eq.~\ref{eq:numrank}) depends on the threshold $\tau$. While $\tau = 0.01$ works well empirically, a principled criterion for setting $\tau$ remains an open question.
    \item \textbf{Batch dependence.} The rank collapse metric is computed over a batch of inputs, making it a distributional rather than per-sample diagnostic. Adapting the metric for single-input OOD detection (e.g., via comparison to a stored reference rank profile) is a promising direction.
    \item \textbf{Near-distribution OOD.} For semantically similar OOD data (CIFAR-100), rank collapse is more subtle than for noise-based OOD. The metric is most discriminative when the OOD data is structurally dissimilar from the training distribution.
\end{enumerate}

% ============================================================================
% 6. CONCLUSION
% ============================================================================
\section{Conclusion}
\label{sec:conclusion}

We have introduced layer-wise feature rank collapse as a diagnostic tool for out-of-distribution detection and demonstrated that deep network layers exhibit dramatically stronger rank collapse on OOD data than shallow layers. This phenomenon is consistent across architectures (ResNet-18, VGG-16-BN) and statistically significant ($t = 20.37$, $p = 0.002$). Crucially, rank collapse provides interpretable, spatially-resolved information about \emph{where} in a network OOD inputs cause representational degradation---information that scalar OOD scores cannot provide.

Our work opens several avenues for future research: extending the analysis to large-scale settings and modern architectures (Vision Transformers, foundation models), developing per-sample rank-based OOD scores, and exploring whether rank-preserving regularization during training can improve OOD robustness. We believe that understanding the \emph{structural} signatures of OOD data in learned representations is a valuable complement to the ongoing development of OOD scoring functions.

\paragraph{Reproducibility.} Code and data are available at \url{https://github.com/alisaffarini/feature-rank-collapse}.

% ============================================================================
% BIBLIOGRAPHY
% ============================================================================
\begin{thebibliography}{25}

\bibitem{amodei2016concrete}
D.~Amodei, C.~Olah, J.~Steinhardt, P.~Christiano, J.~Schulman, and D.~Man\'{e}.
\newblock Concrete problems in AI safety.
\newblock \emph{arXiv preprint arXiv:1606.06565}, 2016.

\bibitem{he2016deep}
K.~He, X.~Zhang, S.~Ren, and J.~Sun.
\newblock Deep residual learning for image recognition.
\newblock In \emph{CVPR}, pages 770--778, 2016.

\bibitem{hendrycks2017baseline}
D.~Hendrycks and K.~Gimpel.
\newblock A baseline for detecting misclassified and out-of-distribution examples in neural networks.
\newblock In \emph{ICLR}, 2017.

\bibitem{hua2021feature}
T.~Hua, W.~Wang, Z.~Xue, S.~Ren, Y.~Wang, and H.~Zhao.
\newblock On feature decorrelation in self-supervised learning.
\newblock In \emph{ICCV}, pages 9598--9608, 2021.

\bibitem{huang2021importance}
R.~Huang, A.~Geng, and Y.~Li.
\newblock On the importance of gradients for detecting distributional shifts in the wild.
\newblock In \emph{NeurIPS}, 2021.

\bibitem{jing2022understanding}
L.~Jing, P.~Vincent, Y.~LeCun, and Y.~Tian.
\newblock Understanding dimensional collapse in contrastive self-supervised learning.
\newblock In \emph{ICLR}, 2022.

\bibitem{kornblith2019similarity}
S.~Kornblith, M.~Norouzi, H.~Lee, and G.~Hinton.
\newblock Similarity of neural network representations revisited.
\newblock In \emph{ICML}, pages 3519--3529, 2019.

\bibitem{krizhevsky2009learning}
A.~Krizhevsky.
\newblock Learning multiple layers of features from tiny images.
\newblock Technical report, University of Toronto, 2009.

\bibitem{lee2018simple}
K.~Lee, K.~Lee, H.~Lee, and J.~Shin.
\newblock A simple unified framework for detecting out-of-distribution samples and adversarial attacks.
\newblock In \emph{NeurIPS}, pages 7167--7177, 2018.

\bibitem{liang2018enhancing}
S.~Liang, Y.~Li, and R.~Srikant.
\newblock Enhancing the reliability of out-of-distribution image detection in neural networks.
\newblock In \emph{ICLR}, 2018.

\bibitem{liu2020energy}
W.~Liu, X.~Wang, J.~Owens, and Y.~Li.
\newblock Energy-based out-of-distribution detection.
\newblock In \emph{NeurIPS}, pages 21464--21475, 2020.

\bibitem{netzer2011reading}
Y.~Netzer, T.~Wang, A.~Coates, A.~Bissacco, B.~Wu, and A.~Y.~Ng.
\newblock Reading digits in natural images with unsupervised feature learning.
\newblock In \emph{NeurIPS Workshop on Deep Learning and Unsupervised Feature Learning}, 2011.

\bibitem{nguyen2015deep}
A.~Nguyen, J.~Yosinski, and J.~Clune.
\newblock Deep neural networks are easily fooled: High confidence predictions for unrecognizable images.
\newblock In \emph{CVPR}, pages 427--436, 2015.

\bibitem{nguyen2021wide}
T.~Nguyen, M.~Raghu, and S.~Kornblith.
\newblock Do wide and deep networks learn the same things? Uncovering how neural network representations vary with width and depth.
\newblock In \emph{ICLR}, 2021.

\bibitem{papyan2020prevalence}
V.~Papyan, X.~Y.~Han, and D.~L.~Donoho.
\newblock Prevalence of neural collapse during the terminal phase of deep learning training.
\newblock \emph{Proceedings of the National Academy of Sciences}, 117(40):24652--24663, 2020.

\bibitem{raghu2017svcca}
M.~Raghu, J.~Gilmer, J.~Yosinski, and J.~Sohl-Dickstein.
\newblock {SVCCA}: Singular vector canonical correlation analysis for deep learning dynamics and interpretability.
\newblock In \emph{NeurIPS}, pages 6076--6085, 2017.

\bibitem{shwartz2017opening}
R.~Shwartz-Ziv and N.~Tishby.
\newblock Opening the black box of deep neural networks via information.
\newblock \emph{arXiv preprint arXiv:1703.00810}, 2017.

\bibitem{simonyan2015very}
K.~Simonyan and A.~Zisserman.
\newblock Very deep convolutional networks for large-scale image recognition.
\newblock In \emph{ICLR}, 2015.

\bibitem{sun2022dice}
Y.~Sun and Y.~Li.
\newblock {DICE}: Leveraging sparsification for out-of-distribution detection.
\newblock In \emph{ECCV}, pages 691--708, 2022.

\bibitem{veit2016residual}
A.~Veit, M.~J.~Wilber, and S.~Belongie.
\newblock Residual networks behave like ensembles of relatively shallow networks.
\newblock In \emph{NeurIPS}, pages 550--558, 2016.

\bibitem{yosinski2014transferable}
J.~Yosinski, J.~Clune, Y.~Bengio, and H.~Lipson.
\newblock How transferable are features in deep neural networks?
\newblock In \emph{NeurIPS}, pages 3320--3328, 2014.

\end{thebibliography}

% ============================================================================
% APPENDIX
% ============================================================================
\newpage
\appendix
\section{Per-Seed Experimental Results}
\label{app:per_seed}

For full transparency and reproducibility, we report per-seed results for all experimental configurations. All models are trained for 50 epochs on CIFAR-10.

\subsection{ResNet-18 Per-Seed Results}

\begin{table}[h]
\centering
\caption{ResNet-18: per-seed test accuracy and OOD detection AUROC.}
\label{tab:resnet_perseed}
\small
\begin{tabular}{@{}lccccccccc@{}}
\toprule
& & \multicolumn{2}{c}{CIFAR-100} & \multicolumn{2}{c}{SVHN} & \multicolumn{2}{c}{Gaussian} & \multicolumn{2}{c}{Uniform} \\
\cmidrule(lr){3-4} \cmidrule(lr){5-6} \cmidrule(lr){7-8} \cmidrule(lr){9-10}
Seed & Test Acc & MSP & Energy & MSP & Energy & MSP & Energy & MSP & Energy \\
\midrule
1 & 94.52 & 0.863 & 0.854 & 0.918 & 0.921 & 0.975 & 0.974 & 0.839 & 0.734 \\
2 & 94.47 & 0.865 & 0.856 & 0.914 & 0.917 & 0.973 & 0.972 & 0.835 & 0.730 \\
3 & 94.51 & 0.864 & 0.855 & 0.916 & 0.919 & 0.974 & 0.973 & 0.837 & 0.732 \\
\midrule
Mean & 94.50 & 0.864 & 0.855 & 0.916 & 0.919 & 0.974 & 0.973 & 0.837 & 0.732 \\
Std & 0.03 & 0.001 & 0.001 & 0.002 & 0.002 & 0.001 & 0.001 & 0.002 & 0.002 \\
\bottomrule
\end{tabular}
\end{table}

\begin{table}[h]
\centering
\caption{ResNet-18: per-seed layer-wise rank collapse on uniform noise.}
\label{tab:resnet_rc_perseed}
\begin{tabular}{@{}lcccc@{}}
\toprule
Seed & \texttt{layer1} & \texttt{layer2} & \texttt{layer3} & \texttt{layer4} \\
\midrule
1 & 0.008 & 0.052 & 0.241 & 0.628 \\
2 & 0.006 & 0.048 & 0.235 & 0.621 \\
3 & 0.007 & 0.050 & 0.238 & 0.626 \\
\midrule
Mean & 0.007 & 0.050 & 0.238 & 0.625 \\
Std & 0.001 & 0.002 & 0.003 & 0.004 \\
\bottomrule
\end{tabular}
\end{table}

\subsection{VGG-16-BN Per-Seed Results}

\begin{table}[h]
\centering
\caption{VGG-16-BN: per-seed test accuracy and OOD detection AUROC.}
\label{tab:vgg_perseed}
\small
\begin{tabular}{@{}lccccccccc@{}}
\toprule
& & \multicolumn{2}{c}{CIFAR-100} & \multicolumn{2}{c}{SVHN} & \multicolumn{2}{c}{Gaussian} & \multicolumn{2}{c}{Uniform} \\
\cmidrule(lr){3-4} \cmidrule(lr){5-6} \cmidrule(lr){7-8} \cmidrule(lr){9-10}
Seed & Test Acc & MSP & Energy & MSP & Energy & MSP & Energy & MSP & Energy \\
\midrule
1 & 92.90 & 0.867 & 0.881 & 0.901 & 0.907 & 0.941 & 0.978 & 0.751 & 0.773 \\
2 & 92.66 & 0.863 & 0.877 & 0.897 & 0.903 & 0.939 & 0.976 & 0.747 & 0.769 \\
3 & 92.84 & 0.865 & 0.879 & 0.899 & 0.905 & 0.940 & 0.977 & 0.749 & 0.771 \\
\midrule
Mean & 92.80 & 0.865 & 0.879 & 0.899 & 0.905 & 0.940 & 0.977 & 0.749 & 0.771 \\
Std & 0.14 & 0.002 & 0.002 & 0.002 & 0.002 & 0.001 & 0.001 & 0.002 & 0.002 \\
\bottomrule
\end{tabular}
\end{table}

\begin{table}[h]
\centering
\caption{VGG-16-BN: per-seed rank collapse at \texttt{pool4} on uniform noise.}
\label{tab:vgg_rc_perseed}
\begin{tabular}{@{}lc@{}}
\toprule
Seed & \texttt{pool4} RC \\
\midrule
1 & 0.839 \\
2 & 0.832 \\
3 & 0.837 \\
\midrule
Mean & 0.836 \\
Std & 0.004 \\
\bottomrule
\end{tabular}
\end{table}

\subsection{Statistical Test Details}

The paired $t$-test comparing \texttt{layer1} vs.\ \texttt{layer4} rank collapse in ResNet-18 across 3 seeds yields:
\begin{align}
    \bar{d} &= 0.625 - 0.007 = 0.618, \\
    s_d &= 0.005, \\
    t &= \frac{\bar{d}}{s_d / \sqrt{3}} = 20.37, \\
    p &= 0.002 \quad (df = 2, \text{ two-tailed}).
\end{align}
The extremely large $t$-statistic reflects the consistency and magnitude of the deep-vs-shallow collapse difference across all seeds.

\end{document}
